% from ega4-9-fr.tex
\oldpage[IV-3]{67}

torsion) est encore \emph{localement constructible}. En effet, soit $s\in S$
et soit $s_0=g(s)$~; il suffira (compte tenu de (1.8.2)) de prouver que, pour
que $\mathcal F_s$ soit de torsion (resp. sans torsion), il faut et il suffit
que $(\mathcal F_0)_{s_0}$ le soit. Or,
$\operatorname{Supp}(\mathcal F_s)$ est l'image réciproque de
$\operatorname{Supp}((\mathcal F_0)_{s_0})$ par la projection
$p:X_s\to(X_0)_{s_0}$ (\textbf{I}, 9.1.13)~; comme $p$ est fidèlement plat et
quasi-compact, dire que $\operatorname{Supp}(\mathcal F_s)$ contient un point
maximal de $X_s$ équivaut à dire que
$\operatorname{Supp}((\mathcal F_0)_{s_0})$ contient un point maximal de
$(X_0)_{s_0}$ (1.1.5 et 2.3.4)~; d'où notre assertion en ce qui concerne
l'ensemble $E$ (\textbf{I}, 7.4.6). Si le sous-Module de torsion $\mathcal T$
de $(\mathcal F_0)_{s_0}$ n'est pas nul,
$\mathcal T\otimes_{k(s_0)}k(s)$ (qui est de torsion d'après ce qui précède)
n'est pas nul et s'identifie à un sous-Module de $\mathcal F_s$ (2.2.7), donc
le sous-Module de torsion de $\mathcal F_s$ n'est pas nul. Enfin, si
$(\mathcal F_0)_{s_0}$ est sans torsion, on peut supposer (en considérant un
ouvert affine de $(X_0)_{s_0}$) que $(\mathcal F_0)_{s_0}$ est isomorphe à un
sous-Module d'un $\mathcal O_{(X_0)_{s_0}}^n$, donc $\mathcal F_s$ est
isomorphe à un sous-Module d'un $\mathcal O_{X_s}^n$ (2.2.7), et ceci établit
notre assertion concernant $E'$.
\end{remark}

\subsection{Constructibilité de propriétés topologiques}
\label{subsection:IV.9.5-fr}

\begin{proposition}[9.5.1]
\label{IV.9.5.1-fr}
Soient $f:X\to S$ un morphisme de présentation finie, $Z$ une partie
localement constructible de $X$. Alors l'ensemble $E$ des $s\in S$ tels que
$Z_s\neq\varnothing$ est localement constructible.
\end{proposition}

En effet, on a $E=f(Z)$, et il suffit d'appliquer le th. de Chevalley (1.8.4).

\begin{corollary}[9.5.2]
\label{IV.9.5.2-fr}
Si $Z'$, $Z''$ sont deux parties localement constructibles de $X$, l'ensemble
des $s\in S$ tels que $Z'_s\subset Z''_s$ \textnormal{(}resp.
$Z'_s=Z''_s$\textnormal{)} est localement constructible.
\end{corollary}

En effet, la relation $Z'_s\subset Z''_s$ équivaut à
$(Z'\cap(\complement Z''))_s=\varnothing$ et
$Z'\cap\complement Z''$ est localement constructible.

\begin{proposition}[9.5.3]
\label{IV.9.5.3-fr}
Soient $f:X\to S$ un morphisme de présentation finie, $Z$, $Z'$ deux parties
localement constructibles de $X$ telles que $Z\subset Z'$. Alors l'ensemble
$E$ des $s\in S$ tels que $Z_s$ soit dense dans $Z'_s$ est localement
constructible dans $S$.
\end{proposition}

Il faut vérifier les deux conditions de (9.2.2, (iii)). En ce qui concerne la
première, considérons un préschéma $X$ algébrique sur un corps $k$, deux parties
constructibles $Z$, $Z'$ de $X$ telles que $Z\subset Z'$, et une extension
$k'$ de $k$. Alors la projection canonique $p:X_{(k')}\to X$ est un morphisme
fidèlement plat et quasi-compact, et l'on a donc
$p^{-1}(\overline Z)=\overline{p^{-1}(Z)}$ et
$p^{-1}(\overline{Z'})=\overline{p^{-1}(Z')}$ en vertu de (2.3.10)~; comme
$p$ est surjectif, la relation $\overline Z=\overline{Z'}$ est équivalente à
$\overline{p^{-1}(Z)}=\overline{p^{-1}(Z')}$.

Vérifions maintenant la seconde condition, et supposons donc $S$ affine,
noethérien et intègre, de point générique $\eta$. Distinguons deux cas~:

$1^\circ$ $\eta\in S-E$, autrement dit, $Z_\eta$ n'est pas dense dans
$Z'_\eta$~; il existe donc dans $X$ un ouvert $V$ tel que
$V\cap Z_\eta=\varnothing$ et $V\cap Z'_\eta\neq\varnothing$. Comme $X$ est
noethérien, $V$ est localement constructible, donc il en est de même de
$V\cap Z$, et en vertu de (9.5.1), il y a un voisinage $U$ de $\eta$ dans $S$
tel que pour tout $s\in U$, on ait $(V\cap Z)_s=\varnothing$ et
$(V\cap Z')_s\neq\varnothing$~; cela entraîne que $Z_s$ n'est pas dense dans
$Z'_s$ pour $s\in U$, autrement dit $U\subset S-E$.

